% Transcription — fonds Alexandre Grothendieck, shelfmark 7, batch 4 (pages 61-71).
% Established from the facsimile archives/batches/7/batch-04.pdf, cut from a
% copy of 7.pdf fetched through the project's own relay
% (https://grothendieck-relay.onrender.com/source/7.pdf) rather than from
% grothendieck.umontpellier.fr directly; per the skill transcribe-grothendieck.
% The batch file carries no cover sheet: PDF page k is the archivists' page
% 60+k, and their pencil numbers bottom-left agree throughout. The folder is
% seventy-one pages; this is the fourth and last batch (eleven pages). Batch 3
% (pp. 41-60) was transcribed in parallel; once it existed this file was
% aligned with it: its p. 60 opens 6.1 « Coh. stratifiante » and breaks off
% on « on a une », which p. 61 continues (« faisceaux ... »), and its
% notation for formal completions, X^{(nu+1)}/X, is followed here.
%
% Pass: Opus 5.5 (claude-opus-5-5), 2026-09-26 — first pass, unchecked against the pages by a human.
%
% The folder sits in the inventory group "4-9" « Cristaux »; it has its own
% date in the catalogue, which \dating{} carries verbatim, with no group
% sentence.
%
% Pages skipped: none. Pages summarised: none. All eleven leaves are his own
% manuscript in blue ink, transcribed as far as the hand allows. Page 67
% carries only one struck line (a first start of 6.5, « Cas de S = Spec C »,
% « Soit X un espace analytique ») and a struck sketch; it is transcribed as
% struck. It interrupts the sentence running from p. 66 to p. 68.
%
% His pagination and dates. None on these leaves: no page numbers of his, no
% date; the nearest dates are batch 3's « 6.12.1966 » (p. 56, under
% « Laïus ») and batch 2's « ?5.11.66 » (p. 35). Page 61 opens mid-sentence (« faisceaux ... »),
% continuing p. 60 of
% batch 3. His paragraph numbers run 6.2 (p. 61), 6.3 (p. 64), 6.4 (p. 66),
% 6.5 (p. 68), each with an underlined heading in the left margin, given as
% \marginal{}: « Coh. cristalline », « Applications aux ... canoniques »,
% « Cas de O_{X_strat} et O_{X_cris} », « Cas S = Spec C ». Paragraph 6.1
% (« Coh. stratifiante ») is on p. 60 of batch 3; §6 is « Calculs
% Čechistes » in the plan of p. 59, and « l'isomorphisme trouvé plus haut »
% (p. 61) is the boxed one of p. 60.
%
% What the batch is about. End of the talk notes: computing the
% cohomology of the stratifying site X_strat and of the crystalline site
% X_cris by Čech-Alexander-type cosimplicial complexes. P. 61: the
% cosimplicial sheaf F^* on X_zar, F^{(nu)}(U) = F(U^{nu+1}/U) = lim_i
% F(Delta^{(nu)}_U(i)), H^*(X_strat, F) = HH^*(X_zar, F^*) and its
% hypercohomology spectral sequence (boxed). 6.2 (pp. 61-63): the
% crystalline case, X embedded in Y formally smooth over S; the sheaf
% Y~ = lim Y(i)~ represented by (X,Y) covers the final object, hence the
% spectral sequence with E_2 = H^p(nu -> H^q(Y~^{nu+1}, F)); for F
% quasi-coherent and X affine, H^*(X_cris, F) = H^*(nu -> F(Y^{nu+1}/X))
% (formal completion along X); for X not affine, the sheaves F^{(nu)} on
% X_zar and H^*(X_cris, F) = HH^*(X_zar, F^*). 6.3 (pp. 64-65): invariance
% under a nilpotent immersion X_0 -> X, H^*(X_cris, F) = H^*(X_0cris, F_0);
% for X smooth over S and S_0 -> S nilpotent, H^*(X_cris, O) =
% H^*(X_0cris, O), so the left side does not depend on the lifting; in
% char. 0 it is H^*(X_zar, Omega^._{X/S}), a definition of De Rham
% cohomology for any f : X -> S, and even Rf_*(Omega^.) as an absolute
% crystal on S. 6.4 (p. 66, 68): F = O_strat or O_cris, F^* a cosimplicial
% ring, coboundaries as differential operators of infinite order; analogy
% with Čech-Alexander cochains. 6.5 (pp. 68-70): S = Spec C; via GAGA and
% Mittag-Leffler, H^*(X_strat, O) = H^*(X^an_strat, O) = HH^*(X^an,
% F^{an*}); Conjecture: for every analytic space the complex F^* is a
% resolution of the constant sheaf C; Th.: then for X proper H^*(X_strat,
% O) = H^*(X^an, C); for X in a smooth Y, H^*(X_cris, O) -> H^*(X_strat,
% O) is an isomorphism (lim over the Y(i)). Pb (pp. 70-71): is H^*(X_cris,
% F) -> H^*(X_strat, F) an isomorphism in general, equivalently invariance
% of H^*(X_strat, F) under a nilpotent immersion. Remarques (p. 71): the
% conjecture for non-archimedean complete valued k (Tate's rigid spaces,
% struck; « flasques » analytic spaces).
%
% Notation fixed once for the batch (continuing batches 1-2): \mathbb{H}
% for his double-stroked H (hypercohomology); \mathcal{F}^{(\nu)},
% \mathcal{F}^{*} for his script F (the cosimplicial sheaf) against roman
% F for the sheaf on the site; \widetilde{Y} for his Y with a tilde (the
% sheaf represented); Y^{\nu+1}/X for the formal completion of Y^{nu+1}
% along X, as batch 3 writes X^{(\nu+1)}/X; X_{\mathrm{zar}},
% X_{\mathrm{cris}}, X_{\mathrm{strat}}. The subscript « strat » is often
% abbreviated to a scrawl (« str. », « strif ») and is set strat throughout,
% noted once on p. 61. On p. 69 his script X for an analytic space, drawn
% like an asterisk crossed, is set \mathfrak{X}. His underlined words in
% running text are set in \emph.
%
% Hardest pages and readings to check first: the struck cartouche and the
% second line of 6.2 on p. 61; the insertion over « Supposons que » on
% p. 64 and his margin there; the underlined word before « com. » on p. 65
% (left \ill{}); the tail of 6.4 on p. 66; the struck line and closing
% sentences of p. 69; the « Pb » of p. 70 and the whole Remarques of
% p. 71, with its hatched cancelled block. Apparatus: 49 \ill{} and 27
% \uncertain{}.

\documentclass[11pt,a4paper]{article}
% Le préambule commun à toutes les transcriptions du fonds Grothendieck.
%
% Il tient en une idée : **l'incertitude se compose, elle ne se résout pas.**
% Une transcription qui lisse un mot douteux a détruit la seule chose pour
% laquelle on la faisait — un lecteur doit pouvoir savoir, à chaque endroit,
% ce qui était sur le feuillet et ce qui a été ajouté. Les six macros qui
% suivent sont donc l'appareil critique, et non de la décoration.
%
% Les mêmes commandes sont reconnues par `scripts/render.mjs`, qui produit la
% vue de lecture du site. Une macro ajoutée ici doit l'être aussi là-bas,
% faute de quoi le rendu échouera bruyamment — ce qui est voulu : un rendu
% silencieusement faux est pire qu'un rendu absent.

\ProvidesPackage{grothendieck}[2026/08/08 Transcriptions du fonds Grothendieck]

\usepackage{amsmath,amssymb,amsthm}
\usepackage{mathrsfs}
\usepackage{stmaryrd}
% Les diagrammes commutatifs sont partout dans ce fonds : c'est la langue
% même de la géométrie algébrique de Grothendieck, et une transcription qui
% les rendrait en prose ne transcrirait rien.
\usepackage{tikz-cd}
% Français par défaut — c'est la langue des feuillets — mais l'anglais est
% chargé aussi : la traduction et le résumé sont des documents anglais, et la
% typographie française (espaces avant les deux-points) y serait une faute.
% Ces documents basculent par \selectlanguage{english} en tête de corps.
\usepackage[english,french]{babel}
\usepackage{fontspec}
\usepackage{geometry}
\geometry{a4paper,margin=25mm,marginparwidth=28mm}
\usepackage{marginnote}
\usepackage{xcolor}
% ulem, not soul. `soul`'s \st and \ul are built on a character-by-character
% scan that XeLaTeX's font machinery breaks: the document fails with an
% undefined control sequence rather than degrading. ulem does the same two jobs
% and is Unicode-engine safe. `normalem` stops it hijacking \emph, which the
% transcriptions use for Grothendieck's own emphasis.
\usepackage[normalem]{ulem}
\usepackage{hyperref}
\hypersetup{colorlinks=true,linkcolor=black,urlcolor=black,pdfborder={0 0 0}}

% --- Métadonnées du lot -----------------------------------------------------
% Reprises telles quelles par le script de rendu, qui les affiche en tête de
% la vue de lecture. Elles fixent ce que le fichier prétend couvrir : sans
% elles, une transcription flotte sans point d'attache dans un fonds de seize
% mille feuillets.

\newcommand{\folder}[1]{\gdef\grothFolder{#1}}
\newcommand{\batch}[1]{\gdef\grothBatch{#1}}
\newcommand{\leaves}[2]{\gdef\grothFirst{#1}\gdef\grothLast{#2}}
\newcommand{\foldertitle}[1]{\gdef\grothFoldertitle{#1}}
\gdef\grothFolder{?}\gdef\grothBatch{?}\gdef\grothFirst{?}\gdef\grothLast{?}\gdef\grothFoldertitle{}

% --- Le résumé --------------------------------------------------------------
%
% Il ouvre la lecture modernisée, sous le titre « Résumé », et il est composé
% en petit. Ce n'est pas un ornement : c'est le seul passage du document qui
% ne soit pas des mathématiques, et un lecteur doit pouvoir le reconnaître
% avant de l'avoir lu — puis le sauter s'il connaît déjà le sujet, ou s'y
% arrêter s'il ne le connaît pas. Le filet qui le suit marque la reprise du
% corps.

\newenvironment{resume}
  {\small\setlength{\parskip}{0.45em}}
  {\par\medskip\hrule\medskip}

% --- Mots-clés ---------------------------------------------------------------
%
% En anglais, en clôture du résumé de la lecture modernisée : le vocabulaire
% moderne sous lequel chercher ce que les feuillets construisent. Le manifeste
% du site (`npm run manifest`) les extrait pour étiqueter la cote entière —
% cette ligne est la source unique des tags, il n'y a pas de fichier à côté.
\newcommand{\keywords}[1]{\par\smallskip\noindent
  {\footnotesize\textsc{Keywords} --- \textit{#1}}\par}

% --- L'appareil critique ----------------------------------------------------

% Le numéro de feuillet, en marge. C'est l'ancre qui permet de régler un
% désaccord sur une formule en regardant *un* feuillet plutôt qu'une chemise
% de six cents. Le site s'en sert aussi pour tourner les pages du fac-similé
% quand on fait défiler la transcription.
\newcommand{\leaf}[1]{%
  \marginnote{\footnotesize\color{gray!70}\textsf{#1}}%
  \label{leaf:#1}\ignorespaces}

% Illisible : on ne devine pas. Un mot inventé qui se lit comme les autres est
% le pire résultat possible.
\newcommand{\ill}{\textcolor{red!60!black}{[\,\ldots\,]}}

% Lecture douteuse : le mot est donné, et signalé comme incertain.
\newcommand{\uncertain}[1]{\uline{#1}}

% Ajout de l'éditeur — une lettre manquante, un mot sous-entendu.
\newcommand{\add}[1]{\textcolor{blue!50!black}{[#1]}}

% Biffé par Grothendieck lui-même. Ce qu'il a rayé fait partie du document :
% une rature dit souvent où la pensée a changé de direction.
\newcommand{\struck}[1]{\sout{#1}}

% Note du transcripteur, sur l'état matériel du feuillet.
\newcommand{\note}[1]{\par\smallskip{\small\color{gray!60!black}\textit{[#1]}}\par\smallskip}

% Note marginale de Grothendieck, distincte du corps du texte.
\newcommand{\marginal}[1]{\par\smallskip\noindent\hspace{1em}%
  {\small\color{gray!60!black}#1}\par\smallskip}

% --- Titre ------------------------------------------------------------------

\AtBeginDocument{%
  \begin{center}
    {\large\bfseries Fonds Alexandre Grothendieck --- cote n\textsuperscript{o}~\grothFolder}\\[2pt]
    {\itshape\grothFoldertitle}\\[4pt]
    {\small feuillets \grothFirst--\grothLast\ (lot \grothBatch)}\\[2pt]
    {\footnotesize\color{gray!60!black}Transcription établie d'après le fac-similé numérisé
     par l'Université de Montpellier.\\
     Les lectures incertaines sont soulignées, les illisibles notées [\,\ldots\,],
     les ajouts entre crochets.}
  \end{center}
  \vspace{1em}}

\endinput


\folder{7}
\batch{4}
\pages{61}{71}
\dating{1966-1967}
\watermark{Édition de démonstration}
\foldertitle{Notes laïus novembre-décembre 1967 (Cristaux) : notes manuscrites (1966-1967)}

\begin{document}

\page{61}\note{la page s'ouvre au milieu d'une phrase commencée p.\ 60}

faisceaux
\[
  \mathcal{F}^{(\nu)} : U \longmapsto F(U^{(\nu+1)}/U)
  = \varprojlim_{i} F\bigl(\Delta^{(\nu)}_{U}(i)\bigr)
\]
sur $X_{\mathrm{zar}}$, soit $\mathcal{F}^{(\nu)}$, et que pour $\nu$
variable, ces faisceaux forment un faisceau cosimplicial sur
$X_{\mathrm{zar}}$. Ceci posé, on aura un isom
\[
  H^{*}(X_{\mathrm{strat}}, F) \simeq
  \mathbb{H}^{*}(X_{\mathrm{zar}}, \mathcal{F}^{*}) ,
  \qquad
  \mathcal{F}^{*} = (\nu \mapsto \mathcal{F}^{\nu})
\]
\note{l'indice « strat » est ici récrit sur un autre mot et abrégé ; il
l'est de même dans toute la suite (« str. », « strat »), et on l'écrit
partout $\mathrm{strat}$. L'égalité $\mathcal{F}^{*} = (\nu \mapsto
\mathcal{F}^{\nu})$ est écrite sous $\mathcal{F}^{*}$, reliée par un signe
$=$ vertical}
qui impliquera une suite spectrale d'hypercoh.
\[
  \boxed{\;H^{*}(X_{\mathrm{strat}}, F) \Longleftarrow
  E_{2}^{pq} = H^{p}\bigl(\nu \mapsto H^{q}(X_{\mathrm{zar}},
  \mathcal{F}^{(\nu)})\bigr)\;}
\]
\note{encadré par lui ; devant $\nu$, un signe biffé}
généralisant l'isomorphisme trouvé plus haut.

\struck{Regardons le complexe $\mathcal{F}^{*}$, prenons pour simplifier
$F = \mathcal{O}_{X_{\mathrm{strat}}}$}\note{ces deux lignes sont biffées,
encadrées et hachurées}

\marginal{\emph{Coh.\ cristalline}}
6.2. Lorsque \uncertain{on} travaille en cohomologie cristalline, \ill{}
\ill{} \uncertain{peut} (\uncertain{lorsque} $X$ n'est pas f\add{ormellemen}t
lisse sur $S$) \ill{} \uncertain{définir} que $\widetilde{X}$ comme l'objet
final. Mais supposons que l'on puisse trouver une immersion

\page{62}

\begin{tikzcd}
  X \arrow[r, hook] \arrow[d] & Y \arrow[dl] \\
  S &
\end{tikzcd}

avec $Y$ formellt.\ lisse sur $S$ [\uncertain{toujours} possible
\struck{localement sur $X$}, \ill{} si $S$, $X$ affines, \ill{} si $X$ est
\uncertain{q.-projectif} sur $S$ \struck{\ill{}} \uncertain{alors} \ill{}
\ill{} \ill{}]. Considérons le faisceau sur le site cristallin de $X/S$
« représenté » par $(X, Y)$ ; de façon précise, si $Y(i)$ est le $i$-ème
voisinage inf.\ de $X$ dans $Y$, on trouve que ce faisceau est le faisceau
\[
  \widetilde{Y} = \varinjlim_{i} \widetilde{Y(i)} .
\]
D'autre part l'hypothèse que $Y$ est f\add{ormellemen}t lisse sur $S$
implique que $\widetilde{Y}$ couvre le faisceau final. D'où une
\[
  H^{*}(X_{\mathrm{cris}}, F) \Longleftarrow
  E_{2}^{pq} = H^{p}\bigl(\nu \mapsto H^{q}(\widetilde{Y}^{\nu+1}, F)\bigr)
\]

\page{63}

Supposons de plus les $F_{X'}$ qu.-coh., \add{et simultanément
\uncertain{bien} si $X'' \hookrightarrow X'$ !}\note{ajouté en haut à
droite, relié par un trait} De plus, supposons $U$ dans $X$ affines, donc
les $Y(i)$ affines. \struck{Comme \ill{}} Comme on a alors
\[
  H^{*}(\widetilde{X'}, F) = H^{*}(X'_{\mathrm{zar}}, F_{X'})
\]
on trouve encore $E_{2}^{pq} = 0$ si $q \neq 0$ d'où
\[
  H^{*}(X_{\mathrm{cris}}, F) \simeq
  H^{*}\bigl(\nu \mapsto F(Y^{\nu+1}/X)\bigr)
\]
\note{sous $Y^{\nu+1}/X$, relié par une accolade : « complété formel,
[le] long de $X$ » ; suivent quelques mots griffonnés et biffés, \ill{}}

Lorsque $X$ n'est pas affine, on introduit encore les faisceaux
\[
  \mathcal{F}^{(\nu)} : U \longmapsto \varprojlim F\bigl(U, U(i)^{\nu+1}\bigr)
\]
sur $X_{\mathrm{zar}}$ [$U(i)$ est le $i$-ème voisinage inf.\ de $U$ dans
$Y^{\nu+1}$ ; si $V$ est un ouvert de $Y$ induisant $U$, cela
\uncertain{se réécrit} \ill{}
\[
  \mathcal{F}^{(\nu)}(U) = F\bigl(V^{\nu+1}/U\bigr)
\]
\note{sous $V^{\nu+1}/U$, relié par une accolade : « complété formel le
long de $U$ »}
] ;

Pour $\nu$ variable, on trouve un faisceau cosimplicial, \struck{dit}
noté $\mathcal{F}^{*}$, et on trouve

\page{64}

\[
  H^{*}(X_{\mathrm{cris}}, F) \simeq
  \mathbb{H}^{*}(X_{\mathrm{zar}}, \mathcal{F}^{*})
\]
d'où
\[
  H^{*}(X_{\mathrm{cris}}, F) \Longleftarrow
  H^{p}\bigl(\nu \mapsto H^{q}(X, \mathcal{F}^{(\nu)})\bigr)
\]

\marginal{\emph{Applications aux \ill{} canoniques}}
\emph{6.3. Cas particulier}. Supposons que \struck{\ill{}}\note{au-dessus
du mot biffé, un ajout : « \ill{} \ill{} »} $X_0 \hookrightarrow X$ une
immersion \emph{nilpotente}, \struck{Alors \ill{}} et que $F$ sur
$X_{\mathrm{cris}}$ provienne d'un $F_0$ sur $X_{0\,\mathrm{cris}}$ par
« restriction ». Alors $\forall$ entier $\nu$, les complétés formels
$Y^{\nu+1}/X_0$ et $Y^{\nu+1}/X$ sont identiques. On trouve donc
\[
  \mathcal{F}^{(\nu)} \simeq \mathcal{F}_0^{(\nu)}
\]
\ill{} sur $X_{\mathrm{zar}} = X_{0\,\mathrm{zar}}$, d'où un isom.
\[
  H^{*}(X_{\mathrm{cris}}, F) \xrightarrow{\ \sim\ }
  H^{*}(X_{0\,\mathrm{cris}}, F_0) .
\]
Ceci a été établi lorsqu'on peut plonger $X$ globalement \struck{dans un}
$Y$ f\add{ormellemen}t lisse sur $S$, mais \struck{\ill{} localement}
utilisant un recouvrement de $X$ par des ouverts affines, on voit que l'on
peut se débarrasser de cette hypothèse. [Il faudrait regarder si on peut se
débarrasser de l'hypothèse type quasi-cohérence faite sur $F$ \ldots ?]

\page{65}

Supposons en particulier que $X$ est \emph{lisse} sur $S$, (donc on peut
prendre $Y = X$, mais peu importe \ldots) et $S_0 \hookrightarrow S$ une
imm.\ nilpotente, d'où $X_0 = X \times_S S_0 \to X$ imm.\ nilpotente.
Considérons $X_0$ comme schéma \emph{sur $S$} (pas sur $S_0$) pour définir
$X_{0\,\mathrm{cris}}$. On trouve alors
\[
  H^{*}(X_{\mathrm{cris}}, \mathcal{O}_{X_{\mathrm{cris}}})
  = H^{*}(X_{0\,\mathrm{cris}}, \mathcal{O}_{X_{0\,\mathrm{cris}}})
\]
\emph{Cette relation prouve que le 1er membre ne dépend pas} (à isom
canonique) \emph{de la façon dont on a relevé la situation lisse
$X_0/S_0$ en une situation lisse sur $S$.} Lorsque $S$ est de car.\ nulle,
on trouve que le premier membre est aussi \add{can.\ isom.\ à}
$H^{*}(X_{\mathrm{zar}}, \Omega^{\cdot}_{X/S})$, ce qui nous donne une
définition des \emph{\ill{}} \emph{com.} \emph{de coh.\ de De Rham}
\add{pour un morphisme quelconque $f : X \to S$}\note{ajouté sous la
ligne, relié par un trait ; à la ligne suivante, un second ajout :
« pour une \ill{} \ill{} \struck{\ill{}} »} \struck{\ill{}} base $S$ de
car.\ nulle. [Même mieux, on \struck{\ill{}} \ill{} avoir à définir
$Rf_{*}(\Omega^{\cdot}_{X/S})$ comme un \emph{cristal absolu} sur $S$
\ldots], \add{\uncertain{donc} l'étendre à \emph{tt} voisinage
infinitésimal de $S$, pas seulement ceux qui \uncertain{se} rétractent sur
$S$]}.

\page{66}

\marginal{\emph{Cas de} $\mathcal{O}_{X_{\mathrm{strat}}}$ \emph{et}
$\mathcal{O}_{X_{\mathrm{cris}}}$}
6.4. Cas $F = \mathcal{O}_{X_{\mathrm{strat}}}$ ou
$\mathcal{O}_{X_{\mathrm{cris}}}$. Alors $\mathcal{F}^{*}$ est un
\emph{anneau cosimplicial} \add{sur $X_{\mathrm{zar}}$}. \struck{Dans le
premier}

Dans le cas de 6.2, c'est la limite projective des anneaux cosimpliciaux
$\mathcal{F}(i)^{*}$, où $\mathcal{F}(i)^{*}$ est relatif à $Y(i)$ (NB.
$Y(i)_{\mathrm{zar}} = X_{\mathrm{zar}}$ !). \uncertain{Ce qui} est donc,
pour l'étude de cet anneau, \ill{} essentiellement au \struck{contexte}
\uncertain{stratifiant}.

Chaque $\mathcal{F}^{(\nu)}$ \add{sur $X_{\mathrm{zar}}$} est une
$\mathcal{O}_{X_{\mathrm{zar}}}$-algèbre de $\nu+1$ façons canoniques, qui
viennent des $\nu+1$ projections $\mathrm{pr}_i$ ($1 \leqslant i \leqslant
\nu+1$) de $X^{\nu+1}$ sur $X$ ou de $X^{\nu+1}/X$ sur $X$. Si on
choisit p.ex.\ $\mathrm{pr}_1$, \add{$\mathcal{F}^{(\nu)}$ devient
\uncertain{pro-quasi-cohérent}.} \struck{les} les opérations simpliciales,
et par suite les opérateurs cobords doivent être considérés comme
\emph{op.\ diff.\ d'ordre infini} : si on \uncertain{désigne} les
$\mathcal{F}^{(\nu)}$, on trouve des op.\ diff.\ ordinaires \ldots
\emph{C'est \uncertain{ici} que complexes de cochaînes canoniques}
(\uncertain{plus fin} que pour les complexes de De Rham), \emph{que les
complexes}\note{la phrase se poursuit p.\ 68}

\page{67}

\marginal{\struck{Cas de $S = \operatorname{Spec} \mathbb{C}$ 6.4.}}
\struck{Soit $X$ un espace analytique. \ill{}}\note{seule ligne de la page,
biffée, avec sous la note marginale un petit croquis biffé : un premier
départ de 6.5}

\page{68}

\emph{d'op.\ différentiels s'introduisent dans la théorie de façon
vraiment naturelle}.

Analogie avec les cochaînes de Čech-Alexander \ldots

\marginal{\emph{Cas} $S = \operatorname{Spec} \mathbb{C}$}
6.5 Supposons $S = \operatorname{Spec} \mathbb{C}$. Regardons la suite
spectrale de cohomologie \uncertain{stratifiante}
\[
  H^{*}(X_{\mathrm{strat}}, \mathcal{O}_{X_{\mathrm{strat}}})
  \Longleftarrow E_{2}^{pq} = H^{p}\bigl(\nu \mapsto
  H^{q}(\mathcal{O}_{X^{\nu+1}/X})\bigr)
\]
\note{devant $\mathcal{O}$, un symbole récrit, \ill{} ; sous
$H^{q}(\ldots)$, relié par une accolade : « cohomologie formelle »}

Si on suppose pour simplifier $X$ propre, donc les
$H^{q}\bigl(\Delta^{(\nu)}(i), \mathcal{O}_{\Delta^{(\nu)}(i)}\bigr)$ de
dim finie, alors on trouve par M.L.
\[
  H^{q}\bigl(X^{\nu+1}/X, \mathcal{O}_{X^{\nu+1}/X}\bigr)
  \simeq \varprojlim_{i} H^{q}\bigl(\Delta^{(\nu)}(i),
  \mathcal{O}_{\Delta^{(\nu)}(i)}\bigr)
\]
et on peut aussi en écrire
\[
  E_{2}^{pq} = \varprojlim_{i} H^{p}\Bigl(\nu \mapsto
  H^{q}\bigl(\Delta^{(\nu)}(i), \mathcal{O}_{\Delta^{(\nu)}(i)}\bigr)\Bigr) .
\]
Or, grâce à GAGA, la cohomologie des $\Delta^{(\nu)}(i)$ peut
s'interpréter de façon analytique. On trouve ainsi
\[
  H^{*}(X_{\mathrm{strat}}, \mathcal{O}_{X_{\mathrm{strat}}})
  \simeq H^{*}(X^{\mathrm{an}}_{\mathrm{strat}},
  \mathcal{O}_{X^{\mathrm{an}}_{\mathrm{strat}}})
  \overset{\mathrm{déf}}{=}
  \mathbb{H}^{*}(X^{\mathrm{an}}, \mathcal{F}^{\mathrm{an}\,*}_{X^{\mathrm{an}}})
\]

\page{69}

où $\mathcal{F}^{\mathrm{an}\,*}_{X}$ se \struck{\ill{}} définit à partir
de $X^{\mathrm{an}}$ \struck{par} \add{de} façon analogue que dans le cas
algébrique :
\[
  \mathcal{F}^{\mathrm{an}\,*}_{X^{\mathrm{an}}}
  = \mathcal{F}^{*}_{X^{\mathrm{an}}}
  = (\nu \mapsto \mathcal{F}^{\nu}_{X^{\mathrm{an}}})
\]
\note{la parenthèse $(\nu \mapsto \ldots)$ est lue sous une rature ; un
premier essai est biffé}
\[
  \mathcal{F}^{\nu}_{X^{\mathrm{an}}} : U \longmapsto
  \Gamma\bigl(U^{\nu+1}/U, \mathcal{O}_{U^{\nu+1}/U}\bigr)
\]
(où $U$ parcourt les ouverts de $X^{\mathrm{an}}$).

\emph{Conjecture}. Pour tout espace analytique $\mathfrak{X}$, le complexe
$\mathcal{F}^{*}_{\mathfrak{X}}$ est une \emph{résolution} du faisceau
$\mathbb{C}_{\mathfrak{X}}$ (via l'augmentation évidente).
\note{$\mathfrak{X}$ rend un X cursif qu'il trace comme une étoile barrée}
[Alors pour $\mathfrak{X}$ propre, on a
$H^{*}(\mathfrak{X}_{\mathrm{strat}}, \mathcal{O}_{\mathfrak{X}_{\mathrm{strat}}})
= H^{*}(\mathfrak{X}, \mathbb{C})$]

\struck{\ill{}}\note{une ligne entière raturée en boucles}
Cela semble plausible, en regardant les « fibres formelles » des
faisceaux \ldots

\emph{Th}. Soit $X$ propre sur $\mathbb{C}$. Si la conjecture plus haut
est vraie pour $X^{\mathrm{an}}$, alors on a un isom.\ can.
\[
  H^{*}(X_{\mathrm{strat}}, \mathcal{O}_{X_{\mathrm{strat}}})
  \simeq H^{*}(X^{\mathrm{an}}, \mathbb{C})
\]

\page{70}

Supposons $X$ contenu dans un schéma lisse $Y$.
\struck{Si la conjecture ci-dessus est vraie pour les $Y(i)^{\mathrm{an}}$
($Y(i) = i$-ème voisinage infinitésimal de $X$ dans $Y$)} alors
l'homom.\ \struck{\ill{} \ill{}}
\[
  H^{*}(X_{\mathrm{cris}}, \mathcal{O}_{X_{\mathrm{cris}}}) \longrightarrow
  H^{*}(X_{\mathrm{strat}}, \mathcal{O}_{X_{\mathrm{strat}}})
\]
est également un isom.\note{les deux lignes « Supposons \ldots{} » et
« est également un isom. » sont marquées d'un trait vertical dans la marge
; le passage biffé est en outre encadré et hachuré}

Le premier isom.\ résulte de l'isom.\ de comparaison plus haut, et de
\[
  H^{*}(X^{\mathrm{an}}, \mathcal{F}^{*}_{X^{\mathrm{an}}})
  \simeq H^{*}(X^{\mathrm{an}}, \mathbb{C}) .
\]
Pour le deuxième, le dévissage du n° 2 \ill{} aussitôt
\[
  H^{*}_{\mathrm{cris}}(X, \mathcal{O}_{X_{\mathrm{cris}}})
  = \varprojlim_{i} H^{*}_{\mathrm{cris}}\bigl(Y(i),
  \mathcal{O}_{Y(i)_{\mathrm{cris}}}\bigr)
\]
Or d'après la première partie, les homom.\ de transition sont des isom.,
et on gagne.

Cela donne \ill{} des \uncertain{rapports} à la conjecture de l'exposé
précédent.

\emph{Pb}. Supposons de façon générale $X$ \emph{\uncertain{plat}}
\add{(loc.\ de prés.\ finie)} sur $S$, et $F$ \add{sur $X_{\mathrm{cris}}$}
satisfaisant aux hyp.\ habituelles (p.\ ex.\ provenant d'un
\uncertain{Module} \ill{} \uncertain{stratifié}) \uncertain{A-t-on}
\ill{} \ill{}
\[
  H^{*}(X_{\mathrm{cris}}, F_{X_{\mathrm{cris}}}) \xrightarrow{\ \sim\ }
  H^{*}(X_{\mathrm{strat}}, F\ (\text{dite stratifiée}))\ .
\]

\page{71}

C'est aussi équivalent au pb suivant : si $X$ est comme dessus, et
$X_0 \to X$ une imm.\ nilp.\ \ill{}, \uncertain{$Y$} de prés.\ finie sur
$S$, a-t-on
\[
  H^{*}(X_{\mathrm{strat}}, F) \xrightarrow{\ \sim\ }
  H^{*}(X_{0\,\mathrm{strat}}, F)
\]
[De plus, on peut supposer $S$, $X$, $X_0$ affines].

\emph{Remarques}. Si $S = \operatorname{Spec} k$, $k$ corps valué complet
non archimédien, il y a lieu de \struck{faire} se poser la question de la
validité de la conjecture \add{ci-dessus}, pour les \struck{\ill{}
espaces rigides-analytiques de Tate \ill{} sur $k$, et algébriques, --
mais c'est \ill{} \ill{} \ill{}}\note{passage biffé, encadré et hachuré}
espaces analytiques \struck{\ill{}} « flasques » sur $k$. Mais cela
\uncertain{n'a} \ill{} pas d'indication globale \ill{}
\uncertain{vraisemblable} \ldots\note{fin de la page et du dossier}

\end{document}
